% EM structure and teaching revision 20260928
\subsection{平面波界面的振幅与能流核对}
两种无耗散、非磁性介质折射率为 $n_1,n_2$，平面波正入射，取同一切向电场正向。电、磁场的切向边界条件给出
\[
 E_i+E_r=E_t,\qquad n_1(E_i-E_r)=n_2E_t.
\]
反射波传播方向反转，故磁场方程中的反射项带负号。解得电场振幅系数
$r=(n_1-n_2)/(n_1+n_2)$、$t=2n_1/(n_1+n_2)$。负的 $r$ 表示电场相位反转，不表示能量反向增加。
能流与 $nE^2$ 成正比，所以反射率与透射率为
\[
 \mathcal R=|r|^2,\qquad \mathcal T=\frac{n_2}{n_1}|t|^2,
 \qquad \mathcal R+\mathcal T=1.
\]
例如空气近似 $n_1=1$、另一介质 $n_2=1.5$，有 $r=-0.2$、$t=0.8$，反射能量占 $4\%$，透射占 $96\%$。把 $t^2=0.64$ 当成透射率会误判能量不守恒。
这些公式使用正入射、线性、无耗散且磁导率相同的前提。斜入射须区分偏振；存在吸收时还要计入介质内的能量沉积。

